\paragraph{window-\texorpdfstring{$6$}{6} 微态 Hilbert 空间的可见/隐藏分解与纤维谱几何}
在定理 \ref{thm:xi-foldbin-gauge-representation-distribution-law-m6-spectrum} 的 \(m=6\) 规范群直积分解、定理 \ref{thm:xi-window6-center-exact-splitting-bdry-cyclic} 的 boundary 中心直和因子，以及定理 \ref{thm:xi-window6-cyclic-kernel-bc3-root-datum} 与表 \ref{tab:fold6_b3c3_root_dictionary_b3}、表 \ref{tab:fold6_b3c3_root_dictionary_c3} 的 \(B_3/C_3\) 根字典已经固定之后，还可把 window-\(6\) 的微态层提升为一条严格的有限维表示论--谱几何链：所有根几何只出现于 gauge 固定点商，而隐藏部门完全由各纤维的零和模组成，并在同纤维团图上量子化为三条有限质量通道。

\begin{theorem}[window-\texorpdfstring{$6$}{6} 微态 Hilbert 空间的规范分解]\label{thm:xi-time-part55d-window6-microstate-hilbert-gauge-splitting}
\leanverified{paper\_xi\_time\_part55d\_window6\_microstate\_hilbert\_gauge\_splitting}
沿用定理 \ref{thm:xi-foldbin-gauge-representation-distribution-law-m6-spectrum} 的记号，记
\[
\Omega_6:=\{0,1,\dots,63\},
\qquad
F_w:=\Fold_6^{\mathrm{bin}}{}^{-1}(w),
\qquad
d(w):=|F_w|
\quad (w\in X_6),
\]
并定义
\[
\mathcal H_6:=\ell^2(\Omega_6)
=
\bigoplus_{w\in X_6}\ell^2(F_w).
\]
在
\[
G_6^{\mathrm{bin}}
=
\prod_{w\in X_6}\mathfrak S_{d(w)}
\]
的自然置换表示下，存在规范正交分解
\[
\mathcal H_6=\mathcal H_{\mathrm{vis}}\oplus \mathcal H_{\mathrm{hid}},
\qquad
\mathcal H_{\mathrm{vis}}
:=
\bigoplus_{w\in X_6}\CC \mathbf 1_{F_w}
=
\mathcal H_6^{G_6^{\mathrm{bin}}},
\]
并且
\[
\mathcal H_{\mathrm{hid}}
\cong
\bigoplus_{w:\,d(w)=4}\mathrm{Std}_4^{(w)}
\oplus
\bigoplus_{w:\,d(w)=3}\mathrm{Std}_3^{(w)}
\oplus
\bigoplus_{w:\,d(w)=2}\mathrm{sgn}_2^{(w)}.
\]
因此
\[
\dim \mathcal H_{\mathrm{vis}}=21,
\qquad
\dim \mathcal H_{\mathrm{hid}}=64-21=43,
\]
且
\[
\dim \mathcal H_{\mathrm{hid}}
=
9\cdot 3+4\cdot 2+8\cdot 1
=
27+8+8.
\]
\end{theorem}
\begin{proof}
由于 \(\{F_w:\ w\in X_6\}\) 构成 \(\Omega_6\) 的不交划分，故
\[
\mathcal H_6
=
\bigoplus_{w\in X_6}\ell^2(F_w)
\]
是规范正交直和。对每个 \(w\in X_6\)，都有
\[
\ell^2(F_w)=\CC \mathbf 1_{F_w}\oplus \mathbf 1_{F_w}^{\perp}.
\]
其中 \(\mathfrak S_{d(w)}\) 在 \(\CC \mathbf 1_{F_w}\) 上平凡作用，在 \(\mathbf 1_{F_w}^{\perp}\) 上给出零和超平面表示。于是当 \(d(w)=4\) 时得到三维标准表示 \(\mathrm{Std}_4\)；当 \(d(w)=3\) 时得到二维标准表示 \(\mathrm{Std}_3\)；当 \(d(w)=2\) 时，\(\mathbf 1_{F_w}^{\perp}\) 为一维符号表示 \(\mathrm{sgn}_2\)。

对整个直积群 \(G_6^{\mathrm{bin}}\) 而言，每个因子 \(\mathfrak S_{d(w)}\) 只作用在对应的 \(\ell^2(F_w)\) 上，因此把上述分解逐纤维相加，便得到所示
\[
\mathcal H_6=\mathcal H_{\mathrm{vis}}\oplus \mathcal H_{\mathrm{hid}}
\]
及隐藏部门的不可约分块。

最后，\(\mathcal H_6^{G_6^{\mathrm{bin}}}\) 恰由各纤维上的常数函数组成，故等于 \(\mathcal H_{\mathrm{vis}}\)。维数计算则由推论 \ref{cor:fold-bin-m6-fiber-hist} 的直方图
\[
4^{\times 9},\qquad 3^{\times 4},\qquad 2^{\times 8}
\]
直接给出。证毕。
\end{proof}

\begin{corollary}[window-\texorpdfstring{$6$}{6} 的 \texorpdfstring{$B_3/C_3$}{B3/C3} 根几何只存在于可见固定点商]\label{cor:xi-time-part55d-window6-root-geometry-visible-fixed-quotient}
沿用定理 \ref{thm:xi-time-part55d-window6-microstate-hilbert-gauge-splitting} 的记号，定义
\[
\mathcal H_{\mathrm{vis}}^{\mathrm{cyc}}
:=
\bigoplus_{w\in X_6^{\mathrm{cyc}}}\CC \mathbf 1_{F_w},
\qquad
\mathcal H_{\mathrm{vis}}^{\mathrm{bdry}}
:=
\bigoplus_{w\in X_6^{\mathrm{bdry}}}\CC \mathbf 1_{F_w}.
\]
则
\[
\mathcal H_{\mathrm{vis}}
=
\mathcal H_{\mathrm{vis}}^{\mathrm{cyc}}
\oplus
\mathcal H_{\mathrm{vis}}^{\mathrm{bdry}}
\cong
\CC^{18}\oplus \CC^3.
\]
并且表 \ref{tab:fold6_b3c3_root_dictionary_b3} 与表 \ref{tab:fold6_b3c3_root_dictionary_c3} 的 \(B_3/C_3\) 字典仅定义在 \(\mathcal H_{\mathrm{vis}}^{\mathrm{cyc}}\) 的自然基上。换言之，window-\(6\) 的根系几何并非微态层上的结构，而是 gauge 固定点商中的可见几何。
\end{corollary}
\begin{proof}
由定理 \ref{thm:xi-window6-cyclic-kernel-bc3-root-datum}，
\[
X_6=X_6^{\mathrm{cyc}}\sqcup X_6^{\mathrm{bdry}},
\qquad
|X_6^{\mathrm{cyc}}|=18,
\qquad
|X_6^{\mathrm{bdry}}|=3.
\]
因此 \(\mathcal H_{\mathrm{vis}}\) 按 cyclic/boundary 精确裂解为
\[
\mathcal H_{\mathrm{vis}}^{\mathrm{cyc}}\oplus \mathcal H_{\mathrm{vis}}^{\mathrm{bdry}}
\cong
\CC^{18}\oplus \CC^3.
\]
另一方面，表 \ref{tab:fold6_b3c3_root_dictionary_b3} 与表 \ref{tab:fold6_b3c3_root_dictionary_c3} 只对 \(X_6^{\mathrm{cyc}}\) 的 \(18\) 个稳定词给出确定性根字典，而不涉及 boundary 三词。故根几何完全落在 \(\mathcal H_{\mathrm{vis}}^{\mathrm{cyc}}\) 上。证毕。
\end{proof}

\begin{theorem}[规范作用像代数与交换子代数的双侧刚性]\label{thm:xi-time-part55d-window6-action-image-commutant-rigidity}
设
\[
\rho:\CC[G_6^{\mathrm{bin}}]\longrightarrow \operatorname{End}(\mathcal H_6)
\]
为自然置换表示的线性延拓。则其像代数满足
\[
\rho\!\bigl(\CC[G_6^{\mathrm{bin}}]\bigr)
\cong
\CC
\oplus
M_3(\CC)^{\oplus 9}
\oplus
M_2(\CC)^{\oplus 4}
\oplus
\CC^{\oplus 8},
\]
而其交换子代数满足
\[
\rho\!\bigl(\CC[G_6^{\mathrm{bin}}]\bigr)'
\cong
M_{21}(\CC)
\oplus
\CC^{\oplus 9}
\oplus
\CC^{\oplus 4}
\oplus
\CC^{\oplus 8}
\cong
M_{21}(\CC)\oplus \CC^{\oplus 21}.
\]
特别地，\(T\in \operatorname{End}(\mathcal H_6)\) 为自伴且与 \(G_6^{\mathrm{bin}}\) 作用可交换，当且仅当
\[
T
=
A_{\mathrm{vis}}
\oplus
\bigoplus_{w\in X_6}\lambda_w I_{\mathbf 1_{F_w}^{\perp}},
\qquad
A_{\mathrm{vis}}=A_{\mathrm{vis}}^*\in \operatorname{End}(\mathcal H_{\mathrm{vis}}),
\quad
\lambda_w\in \RR.
\]
因此，一切与 gauge 作用等变的非平凡耦合都被压缩在 \(21\) 维可见商 \(\mathcal H_{\mathrm{vis}}\) 中；隐藏部门上既不存在跨纤维耦合，也不存在可见--隐藏混合耦合。
\end{theorem}
\begin{proof}
由定理 \ref{thm:xi-time-part55d-window6-microstate-hilbert-gauge-splitting}，
\[
\mathcal H_6
\cong
\mathbf 1^{\oplus 21}
\oplus
\bigoplus_{w:\,d(w)=4}\mathrm{Std}_4^{(w)}
\oplus
\bigoplus_{w:\,d(w)=3}\mathrm{Std}_3^{(w)}
\oplus
\bigoplus_{w:\,d(w)=2}\mathrm{sgn}_2^{(w)}.
\]
对每个 \(w\)，隐藏块 \(\mathbf 1_{F_w}^{\perp}\) 上只有第 \(w\) 个因子 \(\mathfrak S_{d(w)}\) 非平凡作用，其余因子都平凡。因此每个隐藏块都是 \(G_6^{\mathrm{bin}}\) 的不可约表示；若 \(w\neq w'\)，则把表示限制到第 \(w\) 个因子即可看出：\(\mathbf 1_{F_w}^{\perp}\) 上该因子非平凡，而 \(\mathbf 1_{F_{w'}}^{\perp}\) 上该因子平凡，故这些隐藏块两两不同构。

另一方面，Burnside 定理给出
\[
\CC[\mathfrak S_4]\big|_{\mathrm{Std}_4}\cong M_3(\CC),
\qquad
\CC[\mathfrak S_3]\big|_{\mathrm{Std}_3}\cong M_2(\CC),
\qquad
\CC[\mathfrak S_2]\big|_{\mathrm{sgn}_2}\cong \CC.
\]
而 \(21\) 条常数线共同构成平凡表示的 \(21\) 重等型分量，群代数在其上只通过标量作用。于是像代数正是
\[
\CC
\oplus
M_3(\CC)^{\oplus 9}
\oplus
M_2(\CC)^{\oplus 4}
\oplus
\CC^{\oplus 8}.
\]

再由 Schur 引理，交换子在平凡表示的 \(21\) 重等型块上给出全部矩阵代数
\[
\operatorname{End}(\mathcal H_{\mathrm{vis}})\cong M_{21}(\CC),
\]
而在每个 multiplicity-one 的隐藏不可约块上只能给出标量。故
\[
\rho\!\bigl(\CC[G_6^{\mathrm{bin}}]\bigr)'
\cong
M_{21}(\CC)\oplus \CC^{\oplus 21}.
\]
最后，自伴条件只要求 \(A_{\mathrm{vis}}\) 自伴且各标量 \(\lambda_w\) 取实数。由不同不可约分量之间不存在非零交织算子，隐藏部门上不可能出现跨纤维项，也不可能与可见部门混合。证毕。
\end{proof}

\begin{theorem}[粗粒化算子的三值奇异谱]\label{thm:xi-time-part55d-window6-coarse-graining-three-singular-values}
\leanverified{paper\_xi\_time\_part55d\_window6\_coarse\_graining\_three\_singular\_values}
在 \(\ell^2(X_6)\) 的标准正交基 \((e_w)_{w\in X_6}\) 下，定义粗粒化算子
\[
B:\mathcal H_6\longrightarrow \ell^2(X_6),
\qquad
(Bf)(w):=\sum_{n\in F_w}f(n),
\]
并记
\[
D:=\operatorname{diag}(d(w))_{w\in X_6}.
\]
则成立：
\begin{enumerate}
  \item
  \[
  B^*B=\bigoplus_{w\in X_6}J_{d(w)},
  \]
  其中 \(J_d\) 为 \(d\times d\) 全 \(1\) 矩阵；
  \item \(B\) 的非零奇异值恰为
  \[
  \{2^{(9)},\ \sqrt3^{(4)},\ \sqrt2^{(8)}\};
  \]
  \item
  \[
  \ker B=\mathcal H_{\mathrm{hid}},
  \qquad
  \operatorname{rank}B=21;
  \]
  \item 若定义
  \[
  \widetilde B:=D^{-1/2}B,
  \]
  则
  \[
  \widetilde B\widetilde B^*=I_{\ell^2(X_6)},
  \qquad
  \widetilde B^*\widetilde B=P_{\mathrm{vis}},
  \]
  其中 \(P_{\mathrm{vis}}\) 为到 \(\mathcal H_{\mathrm{vis}}\) 的正交投影。特别地，\(\ell^2(X_6)\) 与 \(\mathcal H_{\mathrm{vis}}\) 规范等距同构。
\end{enumerate}
\end{theorem}
\begin{proof}
对每个 \(w\in X_6\)，算子 \(B\) 在 \(\ell^2(F_w)\) 上的限制只是求和泛函：
\[
B|_{\ell^2(F_w)}(f)
=
\left(\sum_{n\in F_w}f(n)\right)e_w.
\]
故
\[
B^*e_w=\mathbf 1_{F_w},
\]
从而 \(B^*B\) 在第 \(w\) 个纤维块上正是 \(J_{d(w)}\)，这就得到第 \(1\) 条。

而 \(J_d\) 的特征值为 \(d\)（重数 \(1\)）与 \(0\)（重数 \(d-1\)），所以 \(B\) 的非零奇异值恰为 \(\sqrt{d(w)}\)。由推论 \ref{cor:fold-bin-m6-fiber-hist} 的直方图，\(d(w)=4,3,2\) 的纤维个数分别为 \(9,4,8\)，于是得到第 \(2\) 条。

同样，\(J_d\) 的零空间正是零和超平面 \(\mathbf 1_{F_w}^{\perp}\)，故
\[
\ker B
=
\bigoplus_{w\in X_6}\mathbf 1_{F_w}^{\perp}
=
\mathcal H_{\mathrm{hid}}.
\]
由
\[
\dim \mathcal H_6=64,
\qquad
\dim \ker B=43
\]
立得 \(\operatorname{rank}B=21\)，从而第 \(3\) 条成立。

最后，直接计算得
\[
BB^*=D,
\]
故
\[
\widetilde B\widetilde B^*
=
D^{-1/2}BB^*D^{-1/2}
=
I_{\ell^2(X_6)}.
\]
另一方面，\(\widetilde B^*\widetilde B\) 在第 \(w\) 个纤维块上等于
\[
d(w)^{-1}J_{d(w)},
\]
这正是到 \(\CC \mathbf 1_{F_w}\) 的正交投影。逐块相加便得到
\[
\widetilde B^*\widetilde B=P_{\mathrm{vis}}.
\]
因此 \(\widetilde B\) 把 \(\mathcal H_{\mathrm{vis}}\) 等距识别为 \(\ell^2(X_6)\)。证毕。
\end{proof}

\begin{theorem}[同纤维团图的质量量子化与热核三指数分解]\label{thm:xi-time-part55d-window6-fiber-clique-mass-quantization}
在 \(\Omega_6\) 上定义图 \(\Gamma_6\)：若且唯若 \(u\neq v\) 且
\[
\Fold_6^{\mathrm{bin}}(u)=\Fold_6^{\mathrm{bin}}(v),
\]
则在 \(u,v\) 之间连一条边。记 \(A_{\Gamma_6}\) 为其邻接矩阵，\(L_{\Gamma_6}\) 为其图拉普拉斯。则
\[
\Gamma_6\cong 9K_4\sqcup 4K_3\sqcup 8K_2.
\]
从而
\[
\operatorname{Spec}(A_{\Gamma_6})
=
\{3^{(9)},\ 2^{(4)},\ 1^{(8)},\ (-1)^{(43)}\},
\]
\[
\operatorname{Spec}(L_{\Gamma_6})
=
\{0^{(21)},\ 4^{(27)},\ 3^{(8)},\ 2^{(8)}\}.
\]
若 \(P_0,P_2,P_3,P_4\) 为 \(L_{\Gamma_6}\) 对应于特征值 \(0,2,3,4\) 的谱投影，则
\[
P_0=P_{\mathrm{vis}},
\qquad
\mathcal H_{\mathrm{hid}}
=
P_2\mathcal H_6\oplus P_3\mathcal H_6\oplus P_4\mathcal H_6,
\]
并且
\[
\dim P_2\mathcal H_6=8,
\qquad
\dim P_3\mathcal H_6=8,
\qquad
\dim P_4\mathcal H_6=27.
\]
进一步，对任意 \(t\ge 0\) 都有精确热核分解
\[
e^{-tL_{\Gamma_6}}
=
P_{\mathrm{vis}}
+
e^{-2t}P_2
+
e^{-3t}P_3
+
e^{-4t}P_4.
\]
因此隐藏部门不含零模，而严格量子化为三条有限质量通道 \(2,3,4\)。
\end{theorem}
\begin{proof}
按定义，\(\Gamma_6\) 的每个连通分支恰是某条纤维 \(F_w\) 上的完全图 \(K_{d(w)}\)。由推论 \ref{cor:fold-bin-m6-fiber-hist}，window-\(6\) 的纤维直方图为
\[
4^{\times 9},\qquad 3^{\times 4},\qquad 2^{\times 8},
\]
故
\[
\Gamma_6\cong 9K_4\sqcup 4K_3\sqcup 8K_2.
\]

对完全图 \(K_d\)，邻接谱为
\[
\{(d-1)^{(1)},\ (-1)^{(d-1)}\},
\]
图拉普拉斯谱为
\[
\{0^{(1)},\ d^{(d-1)}\}.
\]
逐分支相加即得所示总谱。

特别地，\(L_{\Gamma_6}\) 的零特征空间正是每个连通分支上的常数函数空间，因此恰为
\[
\bigoplus_{w\in X_6}\CC \mathbf 1_{F_w}
=
\mathcal H_{\mathrm{vis}}.
\]
相应的正特征空间则是各分支上的零和空间，因而等于 \(\mathcal H_{\mathrm{hid}}\)。其中 \(K_2\) 贡献 \(8\) 维特征值 \(2\) 空间，\(K_3\) 贡献 \(4\cdot 2=8\) 维特征值 \(3\) 空间，\(K_4\) 贡献 \(9\cdot 3=27\) 维特征值 \(4\) 空间。热核公式最后由谱定理直接给出。证毕。
\end{proof}

\begin{conjecture}[六顶点/GFF 极限中的可见--隐藏严格分离]\label{conj:xi-time-part55d-window6-sixvertex-gff-visible-hidden-separation}
沿用定理 \ref{thm:xi-time-part55d-window6-fiber-clique-mass-quantization} 的谱投影 \(P_{\mathrm{vis}},P_2,P_3,P_4\)。设某一临界六顶点尺度极限下，高度函数在适当归一化后收敛到标量高斯自由场。则任意随网格尺度 \(\delta\downarrow 0\) 变化的 window-\(6\) 局部观测量提升
\[
\mathcal O_\delta\in \mathcal H_6
\]
都应满足正交分解
\[
\mathcal O_\delta
=
P_{\mathrm{vis}}\mathcal O_\delta
+
P_2\mathcal O_\delta
+
P_3\mathcal O_\delta
+
P_4\mathcal O_\delta,
\]
并且：
\begin{enumerate}
  \item 只有 \(P_{\mathrm{vis}}\mathcal O_\delta\) 可能产生临界长程极限；
  \item \(P_2\mathcal O_\delta\)、\(P_3\mathcal O_\delta\)、\(P_4\mathcal O_\delta\) 全部只贡献局部且紧的有限质量修正；
  \item window-\(6\) 的隐藏部门不可能再生成第二个独立的无质量 Gaussian 场；
  \item 表 \ref{tab:fold6_b3c3_root_dictionary_b3} 与表 \ref{tab:fold6_b3c3_root_dictionary_c3} 所编码的 \(B_3/C_3\) 根系几何，仅是可见固定点商上的组合记账，而不是新的连续自由场来源。
\end{enumerate}
\end{conjecture}

\begin{remark}
该猜想的要点不在于再次排除额外高斯模，而在于把全部连续极限障碍精确压缩为隐藏部门的三条离散质量 \(2,3,4\)。若其成立，则 window-\(6\) 与六顶点/GFF 的接口将被压缩为一个单一命题：临界性只驻留于 \(21\) 维可见商，而全部隐藏模式只能以有限质量方式附着其上。
\end{remark}

\endinput
